\subsection{Proof of Main Theorem}\label{subsec:approximation-hilbert-recognition}
\paraid{p-0fb4af10c68c}
We apply \Cref{thm:torunczyk} to
$\GHSpace$
using the AR property from Part~\ref{part:topology}
and the discrete approximation in \Cref{thm:discrete-approximation}.

\begin{proof}[Proof of \Cref{thm:part-iv-main}]
\paraid{p-a8eece7b38b0}
By \Cref{lem:gh-basics},
$\GHSpace$
is complete and separable,
and by \Cref{thm:part-iii-main},
it is a nonempty AR.
\Cref{thm:discrete-approximation} with
$\ApproximationDomain_\FamilyIndex=\HilbertCube$
for every
$\FamilyIndex\in\NonnegativeIntegers$
establishes the approximation condition in \Cref{thm:torunczyk}.
Thus
\Cref{thm:torunczyk} implies that
$\GHSpace$
is homeomorphic to
$\SeparableHilbertSpace$.
\end{proof}

\paraid{p-f04f1370b85d}
The proof of \Cref{thm:part-iv-main} also applies to the subspace
of spaces of diameter one.

\begin{corollary}\label{cor:unit-diameter-hilbert-space}
\paraid{p-673b86d6e522}
The subspace
\[
 \UnitDiameterGHSpace
 =\{\MetricSpace{\BaseCarrier}{\MetricSymbol}\in\GHSpace
     \mid\diam\MetricSpace{\BaseCarrier}{\MetricSymbol}=1\}
\]
of
$\GHSpace$
is homeomorphic to the real Hilbert space
$\SeparableHilbertSpace$.
\end{corollary}

\begin{proof}
\paraid{p-a048b3017172}
Since
$\UnitDiameterGHSpace$
is a nonempty closed subspace of
$\GHSpace$,
the space
$\UnitDiameterGHSpace$
is also  complete and separable.
Define
$\DiameterNormalization\colon\GHSpace\setminus\{\SingletonSpace\}\to\UnitDiameterGHSpace$
as follows.
For every
$\MetricSpace{\BaseCarrier}{\MetricSymbol}\in\GHSpace\setminus\{\SingletonSpace\}$,
put
\[
 \DiameterNormalization\MetricSpace{\BaseCarrier}{\MetricSymbol}
 =\MetricSpace{\BaseCarrier}
   {\MetricSymbol/\diam\MetricSpace{\BaseCarrier}{\MetricSymbol}}.
\]
By \Cref{lem:scaling-contraction},
this is a continuous retraction.
Namely,
$\UnitDiameterGHSpace$
is a retract of an open subset of the AR
$\GHSpace$
from \Cref{thm:absolute-extensor}.
This means that
$\UnitDiameterGHSpace$ is
also an ANR.

\paraid{p-e1df20d1720f}
Fix a two-point metric space
$\MetricSpace{D}{\delta}$
of diameter one.
Define
$\ContractionHomotopy\colon\UnitDiameterGHSpace\times\UnitInterval\to\UnitDiameterGHSpace$
as follows.
For every
$\MetricSpace{\BaseCarrier}{\MetricSymbol}\in\UnitDiameterGHSpace$
and every
$\HomotopyTime\in\UnitInterval$,
put
\[
 \ContractionHomotopy(\MetricSpace{\BaseCarrier}{\MetricSymbol},\HomotopyTime)
 =\DiameterNormalization\bigl(
    \MetricQuotient{\BaseCarrier}{(1-\HomotopyTime)\MetricSymbol}
    \times\MetricQuotient{D}{\HomotopyTime\delta}\bigr),
\]
using the maximum metric on the product,
whose diameter is
$\max\{1-\HomotopyTime,\HomotopyTime\}>0$.
\Cref{lem:scaling-contraction,lem:metric-products}
show that
$\ContractionHomotopy$
is continuous and contracts
$\UnitDiameterGHSpace$
to
$\MetricSpace{D}{\delta}$.
Hence
$\UnitDiameterGHSpace$
is an AR by \Cref{thm:contractible-ane,thm:retract-extensor}.

\paraid{p-fa3dc2186b0d}
Let
$\OpenCover_{0}$ be
 a family of  open sets of
 $\GHSpace$
 covering
 $\UnitDiameterGHSpace$
and consider
the open cover
$\OpenCover=\OpenCover_{0}\cup\{\GHSpace\setminus\UnitDiameterGHSpace\}$
of
$\GHSpace$.
Apply the construction in the proof of
\Cref{thm:discrete-approximation}
to maps into
$\UnitDiameterGHSpace$
and the cover
$\OpenCover$.
For every
$\MetricSpace{\BaseCarrier}{\MetricSymbol}\in\UnitDiameterGHSpace$
and every integer
$\EquilateralSize\geq2$,
the maps in Step~\ref{step:discrete-products} satisfy
$ \diam\bigl(\ProductApproximation_\EquilateralSize
             \MetricSpace{\BaseCarrier}{\MetricSymbol}\bigr)
 =\max\{1,\CoverScale\MetricSpace{\BaseCarrier}{\MetricSymbol}\}=1$
because
$\CoverScale\leq1/16$.
Thus the recursive choices and the discreteness argument in Steps~\ref{step:discrete-separation}--\ref{step:discrete-discreteness}
remain valid in
$\UnitDiameterGHSpace$.
The arguments on
closeness to the given cover and discreteness
are still valid for  this subspace.
\Cref{thm:torunczyk} now applies as in the proof of
\Cref{thm:part-iv-main}.
\end{proof}

\begin{remark}\label{rem:banach-homeomorphism}
\paraid{p-70d10e1e63b3}
Kadets proved that all separable infinite-dimensional Banach spaces are
homeomorphic \cite[Theorem]{Kadets1967}.
Consequently,
\Cref{thm:part-iv-main} implies that
$\GHSpace$
is homeomorphic to every real separable infinite-dimensional Banach space,
including
$\ContinuousFunctions([0,1],\RealNumbers)$
with the uniform norm and
$\ell^p$
for every
$1\leq p<\infty$.
\end{remark}

\paraid{p-6349dead9826}
The main theorem also implies strong local contractions.

\begin{corollary}\label{cor:strong-local-contractions}
\paraid{p-82c01acb09b2}
Every point
$\MetricSpace{\BaseCarrier }{\MetricSymbol }\in \GHSpace$
has a basis of open neighborhoods which strongly deformation retract
onto
$\{\MetricSpace{\BaseCarrier }{\MetricSymbol }\}$.
\end{corollary}

\begin{remark}\label{rem:nonarchimedean-gh}
\paraid{p-1f95d530606f}
The author determined the topology of the non-Archimedean
Gromov--Hausdorff space by identifying it isometrically with a function space
\cite[Theorems~1.1 and~1.3(3)--(4)]{Ishiki2024NonseparableUrysohn}.
Let
$\UltrametricGHSpace$
be the set of isometry classes of nonempty compact ultrametric spaces.
Its non-Archimedean Gromov--Hausdorff distance
$\NonArchimedeanGHDistance$
is the infimum of the Hausdorff distances between isometric images
in common ultrametric spaces.
Let
$\UrysohnFunctionModel$
consist of all functions
$f\colon[0,\infty)\to\NonnegativeIntegers$
such that
\[
 f(0)=0,
\]
and for every
$\varepsilon>0$,
\[
 \Card\bigl(\{r\geq\varepsilon\mid f(r)\ne0\}\bigr)<\infty.
\]
Equip
$\UrysohnFunctionModel$
with the ultrametric
\[
 \UrysohnFunctionDistance(f,g)
 =
 \begin{cases}
  \max\{r\in[0,\infty)\mid f(r)\ne g(r)\},&f\ne g,\\
  0,&f=g.
 \end{cases}
\]
Then the cited theorems imply that
\[
 (\UltrametricGHSpace,\NonArchimedeanGHDistance)
 \text{ is isometric to }
 (\UrysohnFunctionModel,\UrysohnFunctionDistance).
\]
\end{remark}
